\documentclass[11pt,letterpaper]{article}
\usepackage[margin=0.72in]{geometry}
\usepackage[T1]{fontenc}
\usepackage[utf8]{inputenc}
\usepackage{lmodern}
\usepackage{microtype}
\usepackage{xcolor}
\usepackage{enumitem}
\usepackage{tabularx}
\usepackage{hyperref}

\definecolor{accent}{HTML}{00419B}
\definecolor{muted}{HTML}{555555}

\hypersetup{
  colorlinks=true,
  urlcolor=accent,
  linkcolor=accent
}

\setlength{\parindent}{0pt}
\setlength{\parskip}{0.45em}
\setlist[itemize]{leftmargin=1.25em, topsep=0.2em, itemsep=0.15em}

\newcommand{\MauriName}{Maurizio Ungaro}
\newcommand{\MauriTitle}{Staff Scientist, Jefferson Lab}
\newcommand{\MauriEmail}{\href{mailto:ungaro@jlab.org}{ungaro@jlab.org}}
\newcommand{\MauriHomepage}{\href{https://maureeungaro.github.io/home/}{maureeungaro.github.io/home}}
\newcommand{\MauriGitHub}{\href{https://github.com/maureeungaro}{github.com/maureeungaro}}
\newcommand{\MauriScholar}{\href{https://scholar.google.com/citations?user=zkWYILYAAAAJ\&hl=en}{Google Scholar}}
\newcommand{\MauriInspire}{\href{https://inspirehep.net/authors/1322331}{INSPIRE}}
\newcommand{\MauriOrcid}{\href{https://orcid.org/0000-0001-6982-3310}{ORCID 0000-0001-6982-3310}}
\newcommand{\MauriLinkedIn}{\href{https://www.linkedin.com/in/maurizio-ungaro-37992062/}{LinkedIn}}
\newcommand{\MauriJLab}{\href{https://www.jlab.org}{Jefferson Lab}}
\newcommand{\MauriHallB}{\href{https://www.jlab.org/physics/hall-b}{Jefferson Lab Hall-B}}
\newcommand{\MauriRPI}{\href{https://www.rpi.edu}{Rensselaer Polytechnic Institute}}
\newcommand{\MauriUniGe}{\href{https://unige.it/}{Universita degli Studi di Genova}}
\newcommand{\MauriGEMC}{\href{https://gemc.github.io/home/}{GEMC}}
\newcommand{\MauriGeantFourJLab}{\href{https://jeffersonlab.github.io/g4home/}{Geant4 at JLab}}
\newcommand{\MauriClasTwelve}{\href{https://github.com/gemc/clas12Tags}{CLAS12 Simulations}}
\newcommand{\MauriOSG}{\href{https://maureeungaro.github.io/home/osg/osg}{CLAS12 on OSG}}
\newcommand{\MauriLTCC}{\href{https://maureeungaro.github.io/home/ltcc/ltcc}{Low Threshold Cherenkov Counter}}
\newcommand{\MauriSPringEight}{\href{https://www.spring8.or.jp/en/}{SPring-8}}
\newcommand{\MauriHPS}{\href{https://www.jlab.org/physics/hall-b/other}{Heavy Photon Search}}

\newcommand{\DocHeader}[1]{
  {\LARGE\bfseries \MauriName}\hfill {\color{muted}#1}\\
  Staff Scientist, \MauriJLab{}\\
  \MauriEmail{} \quad
  \MauriHomepage{} \quad
  \MauriGitHub{}\\
  \MauriScholar{} \quad
  \MauriInspire{} \quad
  \MauriOrcid{}\\
  \vspace{0.4em}
  \hrule
  \vspace{0.8em}
}

\newcommand{\RoleEntry}[4]{
  \textbf{#1}\hfill {\color{muted}#2}\\
  #3\\
  #4
}

\newcommand{\ProjectEntry}[3]{
  \textbf{#1}: #2\\
  {\small\href{#3}{#3}}
}


\begin{document}
\DocHeader{Application Profile}

\section*{Positioning Statement}
I am a nuclear physicist and scientific software developer whose work connects
physics analysis, detector systems, and production-scale simulation
infrastructure. I am particularly interested in roles that value both scientific
judgment and practical software delivery: simulation frameworks, detector
operations, computational science, research software engineering, and technical
leadership.

\section*{Value for Future Roles}
\begin{itemize}
  \item I can bridge domain science and production software: Geant4, \MauriGEMC{}, \MauriClasTwelve{} geometry, digitization, and analysis workflows.
  \item I have experience supporting users through tutorials, examples, technical notes, and operational guidance.
  \item I understand detector operations and calibration through CLAS12 \MauriLTCC{} work, not only simulation abstractions.
  \item I have delivered infrastructure for distributed simulation production on HTCondor and Open Science Grid resources.
  \item I can communicate to multiple audiences: collaborators, students, detector experts, software users, and review committees.
\end{itemize}

\section*{Application Themes}
\textbf{Scientific computing and research software engineering}: emphasize \MauriGEMC{},
Geant4, Python/C++, reproducibility, CI, Docker, and user support.

\textbf{Detector and experimental physics}: emphasize \MauriLTCC{}, \MauriClasTwelve{}, \MauriHallB{},
calibration, detector performance, and nucleon-structure analyses.

\textbf{Collaboration breadth}: where relevant, mention prior work connected to
\MauriSPringEight{} and dark-sector searches through \MauriHPS{}.

\textbf{Leadership and user enablement}: emphasize tutorials, technical notes,
simulation releases, OSG workflows, and cross-collaboration support.

\textbf{Teaching and mentoring}: emphasize undergraduate laboratory instruction,
student supervision at Jefferson Lab, and practical mentoring in data analysis,
hardware calibration, detector simulation, and scientific computing.

\section*{Reusable Short Cover Paragraph}
My work at \MauriJLab{} combines nuclear physics, detector systems, and
scientific software. I develop and support Geant4/\MauriGEMC{} simulation workflows,
maintain CLAS12 production tools, contribute to detector operation and
calibration, and write documentation and tutorials that help users apply complex
simulation tools productively. This combination of physics depth, software
delivery, and user support is the foundation I would bring to a new role.

\section*{Reusable Teaching Paragraph}
My teaching and mentoring experience comes from both formal undergraduate
laboratory instruction and day-to-day research supervision. I taught PHYS 105
and PHYS 152 laboratories at Christopher Newport University, served as a teaching
assistant for Physics I and Electromagnetic Theory at Rensselaer Polytechnic
Institute, and have mentored students at Jefferson Lab on data analysis,
hardware calibration, detector simulation, and physics analysis. My teaching
style emphasizes practical problem solving, experimental judgment, and modern
computing tools that help students connect physical concepts to measurable data.

\section*{Profiles}
\MauriHomepage{} \quad \MauriGitHub{} \quad \MauriScholar{} \quad \MauriInspire{} \quad \MauriOrcid{}

\end{document}
